\documentclass[12pt]{article}
\usepackage[utf8]{inputenc}
\usepackage[T1]{fontenc}
\usepackage[margin=1in]{geometry}
\usepackage{setspace}
\usepackage{xcolor}
\usepackage[colorlinks=true, urlcolor=blue!60!black, linkcolor=blue!60!black]{hyperref}
\setlength{\parindent}{0em}
\setlength{\parskip}{0.6em}
\renewcommand{\baselinestretch}{1.0}
\begin{document}
\begin{center}
{\LARGE \textit{Duncan v.\ Louisiana}, 391 U.S.\ 145 (1968)\par}
\vspace{0.5em}
{\large Case Brief --- Two-Page Edition\par}
\end{center}

\section{Memory Jogger}
The Sixth Amendment jury right applies to the states; a crime punishable by two years is serious and requires a jury (\href{https://supreme.justia.com/cases/federal/us/391/145/}{391 U.S. 145}).

\section{Detailed Case Facts}
Gary Duncan, a nineteen-year-old Black man, stopped to urge his cousins, who had recently integrated a white school, away from an encounter with white boys in Plaquemines Parish, Louisiana. Accused of slapping one boy's elbow, he was charged with simple battery, punishable by up to two years. Louisiana allowed juries only where capital punishment or hard labor was possible, so a judge tried him alone and sentenced him to 60 days and a \$150 fine.

\section{Procedural History}
The Louisiana Supreme Court denied review; the U.S. Supreme Court reversed and remanded.

\section{Judicial Votes}
7--2. White wrote for Warren, Black, Douglas, Brennan, Fortas, and Marshall; Black (with Douglas) and Fortas concurred; Harlan dissented, joined by Stewart.

\section{Holding}
Due process guarantees a jury in every state prosecution that would fall within the Sixth Amendment if tried federally. Petty offenses are exempt, but the maximum authorized penalty is the chief measure of seriousness, and two years is serious.

\section{Analysis of Opinions}
White asked whether the jury is fundamental in the American system rather than in any conceivable fair system, traced its history from Magna Carta to every state constitution, and stressed its role as a check on government oppression. Black reaffirmed total incorporation. Fortas rejected jot-for-jot incorporation of ancillary rules like unanimity. Harlan called selective incorporation illogical, defended bench trials as fair, and warned of a ``constitutional straitjacket'' on the states.

\section{Examples: Future Applications}
\textbf{Same-Side:} shoplifting punishable by twelve months requires a jury, whatever sentence the prosecutor seeks. \textbf{Opposite-Side:} reckless driving punishable by 90 days is petty. \textbf{Fence-Sitter:} four 90-day counts stacked to 360 days of exposure.

\section{Critique}
Critics note that selective incorporation lacks a principle for which rights are in, and that \textit{Williams} and \textit{Apodaca} diluted the jury until \textit{Ramos}. Civil rights scholars stress the limits of a jury right without representative juries; originalists prefer the Privileges or Immunities Clause.

\section{Key Quotations}
\begin{quote}
``Because we believe that trial by jury in criminal cases is fundamental to the American scheme of justice, we hold that the Fourteenth Amendment guarantees a right of jury trial in all criminal cases which --- were they to be tried in a federal court --- would come within the Sixth Amendment's guarantee.''
\end{quote}
\noindent --- Justice White, Opinion of the Court, 391 U.S. at 149.

\end{document}
